%%%%%%%%%%%%%%%%%%%%%%%%%%%%%%%%%%%%%%%%%%%%%%%%%%%%%%%%%%%%%%%%%%%%%%
% Lớp 4 | Chủ đề 1. Ôn tập và bổ sung | Ôn tập các phép tính trong phạm vi 100 000
% Dạng 4D1?2-2: Phép tính bị che: điền chữ số còn thiếu
%%%%%%%%%%%%%%%%%%%%%%%%%%%%%%%%%%%%%%%%%%%%%%%%%%%%%%%%%%%%%%%%%%%%%%


% nguồn: 2026-11-18_lop4_chu-de-1-tu-luan (nhập 18/11/2026)
\begin{bt}%[4D1N2-2]%[ma:d7629dc1]
% nguon: Toán tư duy lớp 4 – Bài 2 SGK Toán 4 KNTT
	Tìm các chữ số $a$, $b$ trong phép cộng: $\overline{4a6} + \overline{25b} = 712$.
	\loigiai{
		Hàng đơn vị: $6 + b$ có tận cùng là $2$ nên $6 + b = 12$, suy ra $b = 6$, nhớ $1$.\\
		Hàng chục: $a + 5 + 1$ có tận cùng là $1$ nên $a + 6 = 11$, suy ra $a = 5$, nhớ $1$.\\
		Hàng trăm: $4 + 2 + 1 = 7$, đúng.\\
		Vậy $a = 5$, $b = 6$. Thử lại: $456 + 256 = 712$.
	}
\end{bt}

% nguồn: 2026-11-18_lop4_chu-de-1-tu-luan (nhập 18/11/2026)
\begin{bt}%[4D1H2-2]%[ma:d45b4639]
% nguon: Toán tư duy lớp 4 – Bài 2 SGK Toán 4 KNTT
	Tìm các chữ số $a$, $b$ trong phép trừ: $625 - \overline{1a8} = \overline{b57}$.
	\loigiai{
		Hàng đơn vị: $5$ không trừ được $8$, ta lấy $15 - 8 = 7$, nhớ $1$ sang hàng chục.\\
		Hàng chục: $2$ trừ (nhớ $1$ và $a$) phải ra $5$. Vì $2 < 5$ nên phải mượn: $12 - 1 - a = 5$, suy ra $a = 6$, nhớ $1$ sang hàng trăm.\\
		Hàng trăm: $6 - 1 - 1 = 4$, suy ra $b = 4$.\\
		Vậy $a = 6$, $b = 4$. Thử lại: $625 - 168 = 457$.
	}
\end{bt}

% nguồn: 2026-11-18_lop4_chu-de-1-tu-luan (nhập 18/11/2026)
\begin{bt}%[4D1V2-2]%[ma:f655ae30]
% nguon: Toán tư duy lớp 4 – Bài 2 SGK Toán 4 KNTT
	Tìm các chữ số $a$, $b$, $c$ ($a$ khác $0$) biết: $\overline{abc} + \overline{ab} + a = 271$.
	\loigiai{
		Ta có $\overline{abc} = a \times 100 + b \times 10 + c$ và $\overline{ab} = a \times 10 + b$.\\
		Cộng lại: $a \times 111 + b \times 11 + c = 271$.\\
		Nếu $a = 1$ thì $b \times 11 + c = 160$; nhưng lớn nhất $9 \times 11 + 9 = 108$, không được.\\
		Nếu $a \ge 3$ thì $a \times 111 \ge 333 > 271$, không được. Vậy $a = 2$.\\
		Khi đó $b \times 11 + c = 271 - 222 = 49$. Vì $c \le 9$ nên $b = 4$, $c = 49 - 44 = 5$.\\
		Vậy $a = 2$, $b = 4$, $c = 5$. Thử lại: $245 + 24 + 2 = 271$.
	}
\end{bt}

% nguồn: 2026-11-18_lop4_chu-de-1-tu-luan-bo-sung (nhập 18/11/2026)
\begin{bt}%[4D1V2-2]%[ma:4a5a6866]
% nguon: Toán tư duy lớp 4 – Bài 2 SGK Toán 4 KNTT
	Tìm số có hai chữ số $\overline{ab}$, biết: $\overline{ab} \times 5 = \overline{1b5}$.
	\loigiai{
		$\overline{1b5}$ nằm trong khoảng từ $100$ đến $199$, nên $\overline{ab}$ từ $20$ đến $39$ ($20 \times 5 = 100$, $40 \times 5 = 200$).\\
		$\overline{1b5}$ có tận cùng là $5$ nên $\overline{ab}$ là số lẻ: $21, 23, \dots, 39$.\\
		Ta cần chữ số hàng chục của tích đúng bằng chữ số $b$ (hàng đơn vị của $\overline{ab}$):\\
		$21 \times 5 = 105$, $23 \times 5 = 115$, $25 \times 5 = 125$, $27 \times 5 = 135$, $29 \times 5 = 145$: không số nào được.\\
		$31 \times 5 = 155$, $33 \times 5 = 165$, $35 \times 5 = 175$, $37 \times 5 = 185$: không được; $39 \times 5 = 195$: được!\\
		Vậy $\overline{ab} = 39$.
	}
\end{bt}
